\documentclass[10pt]{article}

\usepackage[a4paper,margin=1in]{geometry}
\usepackage{amsmath,amssymb}
\usepackage{booktabs}
\usepackage{enumitem}
\usepackage{titlesec}
\usepackage[colorlinks=true,linkcolor=black,urlcolor=blue]{hyperref}
\setlist{leftmargin=*,itemsep=0.16em,topsep=0.24em}
\setlength{\parindent}{0pt}
\setlength{\parskip}{0.3em}
\linespread{0.97}
\titlespacing*{\section}{0pt}{0.75em}{0.32em}
\titlespacing*{\subsection}{0pt}{0.52em}{0.18em}

\title{\textbf{Poster Presentation Script}\\
Influence Maximization When Communication Is the Bottleneck}
\author{Yixin Chen \and Bichi Zhang \and Yaoqin Ye}
\date{CS240: Algorithm Design and Analysis, Spring 2026\\June 18, 2026}

\begin{document}
\maketitle

\section*{How to Use This Script}

This script is written for a poster-side presentation, not a slide talk. The
recommended style is to point to the poster from left to right:
\begin{enumerate}
  \item start from the research question at the top;
  \item explain the optimization target, diffusion model, and submodularity on
  the left;
  \item walk through the algorithmic story and distributed protocol in the
  center;
  \item end with the large-scale evidence, communication tradeoff, and
  takeaways on the right.
\end{enumerate}

The main version below is designed for about 4--5 minutes. A shorter version is
included for a quick 90-second explanation.

\section{Main Poster Talk: 4--5 Minutes}

\subsection*{Opening}

Hello, we are presenting our CS240 course project, \emph{Influence Maximization
When Communication Is the Bottleneck}. The broad problem is influence
maximization: given a network, choose a small set of initial seed nodes so that
information or behavior spreads to as many nodes as possible.

Our specific question is shown at the top of the poster: how much influence
quality can we keep when expensive centralized seed selection is replaced by
cached greedy evaluation or communication-limited local candidates?

So the project has two connected parts. First, we reproduce the classical
algorithmic story: influence maximization under the Independent Cascade model
is a monotone submodular maximization problem, so greedy has a principled
approximation guarantee. Second, we ask what changes when the selection process
is made more scalable by CELF lazy evaluation or by GreeDI-style distributed
candidate sharing.

\subsection*{Problem and Model}

On the left side of the poster, we define the optimization target. Given a
directed graph \(G=(V,E,p)\) and a seed budget \(k\), we choose
\[
S\subseteq V,\qquad |S|\le k,
\]
to maximize the expected spread
\[
\sigma(S).
\]
In our implementation, \(\sigma(S)\) is estimated by Monte Carlo simulation
under the Independent Cascade model.

The Independent Cascade model is simple but expressive. When a node first
becomes active, it gets one chance to activate each inactive out-neighbor.
Each edge attempt succeeds independently with probability \(p_{uv}\). For our
experiments, we use the weighted-cascade rule,
\[
p_{uv}=\frac{1}{\mathrm{indegree}(v)}.
\]
This means that nodes with many incoming neighbors are harder to activate from
any single incoming edge.

The key theoretical point is submodularity. The expected influence function is
monotone and submodular, meaning that the marginal gain of adding a new seed
decreases as the current seed set grows:
\[
\sigma(A\cup\{v\})-\sigma(A)
\ge
\sigma(B\cup\{v\})-\sigma(B)
\quad\text{for } A\subseteq B.
\]
This is why greedy is not just a heuristic here. For the exact objective,
greedy has the classical \((1-1/e)\) guarantee. The practical bottleneck is
that each marginal gain estimate requires many diffusion simulations.

\subsection*{Algorithms}

The center of the poster shows the algorithmic story.

We start with simple structural baselines: random selection, weighted
out-degree, and PageRank. These are cheap and help us understand how much of
the performance comes from graph structure alone.

Then we implement centralized greedy. At each round, greedy evaluates the
estimated spread of adding each remaining candidate node, and chooses the node
with the largest marginal gain. This matches the submodular optimization
theory, but it is expensive because it repeatedly calls the spread estimator.

Next is CELF, or Cost-Effective Lazy Forward. CELF uses a priority queue of
cached marginal gains. Because submodularity tells us marginal gains can only
go down as the seed set grows, old marginal gains are valid upper bounds. CELF
therefore recomputes a node only when necessary. It targets the same greedy
objective, but reduces redundant spread evaluations.

Finally, we study a GreeDI-style distributed method. The node set is randomly
partitioned into \(m\) parts. Each local worker runs greedy over its local
candidate nodes and sends at most \(q\) candidates to the coordinator. The
coordinator takes the union of local candidates and runs a final greedy
selection over that smaller pool.

This gives us the main distributed tradeoff: smaller \(q\) means lower
communication, but it may miss useful candidates; larger \(q\) sends more
candidates and gives the coordinator more choice, but it increases both
communication and computation.

\subsection*{Experimental Setup}

For experiments, we use synthetic graph families and one small real graph. The
large preset uses Erdős--Rényi, Barabási--Albert, and Watts--Strogatz graphs
with \(n\in\{50,100,200\}\), plus the 34-node Karate Club graph. The seed
budget is \(k=10\) for the large synthetic graphs and \(k=5\) for Karate.

For distributed experiments, we sweep
\[
m\in\{2,4,8\},\qquad q\in\{k,2k,5k\}.
\]
To make comparisons fair, after each algorithm selects its seed set, we
re-evaluate the final seed set using a common evaluation seed for that graph.
So the reported spread is not just the internal value used during selection.

\subsection*{Results}

The right side of the poster summarizes the main results.

First, CELF is the clearest algorithmic improvement. On the large preset,
centralized greedy uses about 1026.5 spread evaluations on average, while CELF
uses 266.2. That is about a 74.1\% reduction. Runtime also drops from about
4.67 seconds to 0.77 seconds on average, an 83.5\% reduction. At the same time,
the average spread stays very close to greedy: 40.05 for CELF versus 40.43 for
greedy.

Second, weighted degree is surprisingly strong. Its average spread is 43.71,
which is above our greedy result under fair re-evaluation. This does not
contradict the greedy approximation theorem. The theorem is about exact greedy
on the exact submodular objective. Our implementation uses Monte Carlo
estimates, and the weighted-cascade probability rule naturally favors
structurally central nodes. So this is a useful empirical lesson: on small
course-scale graphs, graph structure and estimation noise both matter.

Third, the distributed method is competitive but not monotone in the
communication budget. The best average distributed setting in the large preset
is \(m=2, q=k\), which is essentially tied with centralized greedy after fair
re-evaluation. However, increasing \(q\) does not automatically improve
quality in this run. It clearly increases transmitted candidates and spread
evaluations, but the quality still depends on partitions and Monte Carlo
noise.

\subsection*{Conclusion}

The main takeaway is that influence maximization gives a clean bridge from
classical algorithms to modern network applications. Submodularity gives a
principled greedy baseline. CELF turns the same theory into a practical speedup
by avoiding unnecessary evaluations. Distributed candidate sharing changes the
cost model: it can reduce global search, but it introduces communication
limits, partition sensitivity, and more dependence on estimation stability.

So our project is not a production-scale influence maximization system. It is
a course-scale reproduction of the main algorithmic mechanisms: greedy
submodular optimization, lazy evaluation, and communication-constrained
distributed selection.

\section{Short Poster Talk: 90 Seconds}

Our project studies influence maximization: choose \(k\) seed nodes in a
network to maximize expected diffusion spread. Under the Independent Cascade
model, the expected spread function is monotone and submodular, so greedy has
the classical \((1-1/e)\) guarantee for the exact objective.

The practical problem is that greedy is expensive. Every round tries many
candidate nodes, and each candidate requires Monte Carlo spread estimation. So
we compare three levels of methods: simple structural baselines such as random,
weighted degree, and PageRank; centralized greedy and CELF lazy greedy; and a
GreeDI-style distributed method where local workers send only a limited number
of candidates to a coordinator.

Our clearest result is that CELF keeps almost the same spread as greedy but
uses many fewer spread evaluations. On the large preset, evaluations drop from
1026.5 for greedy to 266.2 for CELF, and runtime drops from 4.67 seconds to
0.77 seconds on average.

For distributed candidate sharing, the main message is a tradeoff. Small
candidate pools can be competitive, but larger \(q\) increases communication
and computation and does not automatically improve quality under noisy Monte
Carlo estimation.

Overall, the project shows how submodular optimization explains the greedy
baseline, how CELF accelerates it, and how communication limits complicate
influence maximization in distributed settings.

\section{Common Questions and Suggested Answers}

\subsection*{Q1. Why is influence maximization NP-hard?}

Kempe, Kleinberg, and Tardos showed that influence maximization is NP-hard
under standard diffusion models such as Independent Cascade and Linear
Threshold. Intuitively, the seed set has to solve a combinatorial coverage-like
problem over many possible diffusion outcomes. This is why we use greedy
approximation rather than exact optimization.

\subsection*{Q2. What exactly is submodularity, and why does it matter here?}

Submodularity means diminishing returns. If \(A\subseteq B\), then adding a
candidate node \(v\) to the smaller set \(A\) gives at least as much marginal
gain as adding it to the larger set \(B\). This matters because monotone
submodular maximization under a cardinality constraint has a \((1-1/e)\)
greedy approximation guarantee.

\subsection*{Q3. Does CELF change the greedy solution?}

Conceptually, CELF is an acceleration of greedy, not a different objective. It
uses cached marginal gains as upper bounds and recomputes only when needed. In
exact deterministic evaluation, it can reproduce the greedy choices with far
fewer evaluations. In our implementation, Monte Carlo randomness means the
selected seed set can differ slightly, but the target remains the same greedy
objective.

\subsection*{Q4. Why does weighted degree sometimes beat greedy in your table?}

This is because our reported spread is a fair re-evaluation of final seed
sets, while selection itself uses noisy Monte Carlo estimates. Also, the
weighted-cascade probability rule makes structurally central nodes very strong.
So this result does not violate the greedy theorem; it shows that on
course-scale graphs, estimation noise and graph structure can make a cheap
baseline look very strong.

\subsection*{Q5. Why use Monte Carlo simulation instead of exact spread?}

Exact expected spread is computationally expensive because it averages over
many possible diffusion outcomes. Monte Carlo simulation is simple and
standard for course-scale experiments: run the diffusion process many times
and average the final active set size.

\subsection*{Q6. Why did you use the weighted-cascade probability rule?}

Weighted cascade, \(p_{uv}=1/\mathrm{indegree}(v)\), is common in influence
maximization experiments. It prevents high-indegree nodes from being too easy
to activate by any one neighbor. It also gives a simple, reproducible way to
assign probabilities to synthetic graphs.

\subsection*{Q7. What is the difference between greedy and PageRank or degree?}

Degree and PageRank are structural heuristics. They rank nodes based on graph
structure without simulating the diffusion objective for every possible seed
set. Greedy directly estimates marginal influence spread at each step, so it
is closer to the formal objective but much more expensive.

\subsection*{Q8. Why does distributed greedy evaluate spread on the full graph
instead of only local subgraphs?}

In our simplified model, each worker has a local candidate list, but influence
can still spread globally. Evaluating on the full graph models the fact that a
locally selected seed may activate nodes outside its partition. This keeps the
experiment focused on candidate communication limits rather than on cutting
off diffusion at partition boundaries.

\subsection*{Q9. Why does increasing \(q\) not always improve distributed
quality?}

Increasing \(q\) gives the coordinator more candidates, but the candidates are
still selected using noisy Monte Carlo estimates. Larger \(q\) also changes
the candidate pool and can introduce more redundant candidates. With small
graphs and limited simulations, quality is not guaranteed to increase
monotonically.

\subsection*{Q10. What are the main limitations of this project?}

The biggest limitations are scale and randomness. We use course-scale graphs
up to 200 synthetic nodes, not million-node real networks. We also use a
modest number of Monte Carlo simulations for runtime reasons. A stronger study
would use larger real datasets, more random seeds, confidence intervals, and a
parallel implementation for the distributed workers.

\subsection*{Q11. Did you reproduce IMM, RIS, PaC-IM, or GreediRIS?}

No. Those are treated as related work and motivation. Our implementation
focuses on the core mechanisms that fit the course scope: IC simulation,
centralized greedy, CELF lazy greedy, and GreeDI-style communication-limited
candidate sharing.

\subsection*{Q12. What would be the next extension?}

The most natural extension would be to average over multiple random seeds and
report confidence intervals. Algorithmically, the next step would be an
RR-set-based maximum coverage method, or a truly parallel implementation of
the worker phase to measure real distributed speedup.

\section{Backup Details}

\subsection*{Key Numbers to Remember}

\begin{center}
\begin{tabular}{lrrrr}
\toprule
Method & Spread & Ratio & Runtime & Evaluations \\
\midrule
Greedy & 40.43 & 1.00 & 4.67s & 1026.5 \\
CELF & 40.05 & 0.99 & 0.77s & 266.2 \\
Weighted degree & 43.71 & 1.07 & 0.01s & 1.0 \\
\bottomrule
\end{tabular}
\end{center}

\subsection*{One-Sentence Summary}

Our project shows that submodularity gives influence maximization a principled
greedy baseline, CELF makes that baseline much cheaper, and distributed
candidate sharing introduces a communication-quality tradeoff rather than a
free speedup.

\end{document}
